% AA-70
\chapter[Wzór charakterów Weyla]{Wzór charakterów Weyla}
\label{aa:ch:weyl-character-formula}

Z tożsamości mianownika AA-69 wynika, że iloczyn charakteru i mianownika jest antyniezmienniczy. Nie wynika jednak, że jest pojedynczym alternantem. Wykluczymy niższe alternanty, wykorzystując operator Casimira i ślady na przestrzeniach wagowych.

W rozdziale $K$ jest algebraicznie domkniętym ciałem charakterystyki zero, $\mathfrak g$ niezerową skończenie wymiarową półprostą algebrą Liego nad $K$, $\mathfrak h$ maksymalną podalgebrą toralną, $\Phi^+$ zbiorem korzeni dodatnich względem ustalonej bazy $(\alpha_i)$, $P$ kratą wag, $P_+$ zbiorem wag dominujących, a $W$ grupą Weyla. Używamy $\rho=\frac12\sum_{\alpha\in\Phi^+}\alpha$, $A_\nu$ i $D=A_\rho$ z AA-69. Przez $\kappa$ oznaczamy formę Killinga całej algebry $\mathfrak g$.

\section{Casimir i krotności wag}

Dla $\mu\in\mathfrak h^*$ niech $t_\mu\in\mathfrak h$ będzie określony przez $\kappa(t_\mu,h)=\mu(h)$ dla wszystkich $h\in\mathfrak h$; istnieje jednoznacznie dzięki Twierdzeniu~\ref{aa:tw:maximal-toral-cartan}. Piszemy $(\mu,\nu)=\kappa(t_\mu,t_\nu)$. Na kracie $P$ wartości są wymierne: w bazie $h_i$ macierz $\kappa|_{\mathfrak h}$ ma całkowite wpisy i niezerowy wyznacznik, a współrzędne wag są całkowite. W rzeczywistym modelu AA-60 ta sama forma jest dodatnio określona i niezmiennicza względem $W$ (Twierdzenie~\ref{aa:tw:lie-root-realization}). Zapis $\|\mu\|^2=(\mu,\mu)$ odnosi się do tej formy. Nierówności liczb wymiernych interpretujemy w $\mathbb R$, nie w $K$.

\begin{lemat}
\label{aa:lem:weyl-casimir-eigenvalue}
Przy powyższych założeniach operator Casimira $C_\kappa$ na $L(\lambda)$ dla $\lambda\in P_+$ działa jako skalar
\[
 c_\lambda=(\lambda,\lambda+2\rho)
 =\|\lambda+\rho\|^2-\|\rho\|^2.
\]
\end{lemat}
\begin{proof}
Niech $(H_j)$ i $(H^j)$ będą bazami $\mathfrak h$ sprzężonymi względem $\kappa$. Dla każdego $\alpha\in\Phi^+$ wybieramy $X_\alpha\in\mathfrak g_\alpha$ i $Y_\alpha\in\mathfrak g_{-\alpha}$ tak, że $\kappa(X_\alpha,Y_\alpha)=1$. Jest to możliwe dzięki niezdegenerowanemu parowaniu z Twierdzenia~\ref{aa:tw:root-killing-pairing} i jednowymiarowości przestrzeni korzeniowych z Twierdzenia~\ref{aa:tw:root-dim-reduced}. Niezmienniczość formy daje dla $h\in\mathfrak h$
$\kappa([X_\alpha,Y_\alpha],h)=\kappa(X_\alpha,[Y_\alpha,h])=\alpha(h)$.
Ponieważ nawias należy do $\mathfrak h$, mamy $[X_\alpha,Y_\alpha]=t_\alpha$. Ortogonalność różnych składowych poza parami przeciwległych korzeni daje
\[
 C_\kappa=\sum_j H_j H^j+
 \sum_{\alpha\in\Phi^+}(X_\alpha Y_\alpha+Y_\alpha X_\alpha)
 =\sum_j H_j H^j+t_{2\rho}
       +2\sum_{\alpha\in\Phi^+}Y_\alpha X_\alpha .
\]
Produkty oznaczają złożenia operatorów na module. Na przestrzeni wagi $\mu$ pierwsza suma działa przez $\sum_j\mu(H_j)\mu(H^j)=(\mu,\mu)$, ponieważ $t_\mu=\sum_j\mu(H^j)H_j$. Drugi składnik działa przez $(\mu,2\rho)$. Na wektorze najwyższej wagi $\lambda$ wszystkie $X_\alpha$ działają zerowo: ich obrazy miałyby wagę $\lambda+\alpha$, która nie jest $\preceq\lambda$. Otrzymujemy skalar $c_\lambda$ na tym wektorze. Operator komutuje z działaniem algebry na mocy Lematu~\ref{aa:lem:lie-casimir-identity}. Ponieważ wektor najwyższej wagi generuje $L(\lambda)$, ten sam skalar działa na całym module.
\end{proof}

\begin{lemat}[Rekurencja krotności]
\label{aa:lem:weyl-weight-trace}
Przy powyższych założeniach niech $m_\mu=\dim L(\lambda)_\mu$, także dla wag niewystępujących (wtedy $m_\mu=0$). Dla $\mu\in P$ zachodzi
\[
 \bigl(c_\lambda-(\mu,\mu+2\rho)\bigr)m_\mu
 =2\sum_{\alpha\in\Phi^+}\sum_{j\ge1}
       (\mu+j\alpha,\alpha)m_{\mu+j\alpha}.
\]
Każda suma po $j$ jest skończona po odrzuceniu zerowych składników.
\end{lemat}
\begin{proof}
Dla ustalonego $\alpha\in\Phi^+$ oznaczmy przez $T_\alpha(\mu)$ ślad $Y_\alpha X_\alpha$ na przestrzeni wagi $\mu$. Operatory $X_\alpha$ i $Y_\alpha$ przechodzą odpowiednio z przestrzeni wagi $\mu$ do $\mu+\alpha$ i z powrotem. Dla dowolnych liniowych odwzorowań $A:U\to V$, $B:V\to U$ w skończonych wymiarach mamy $\operatorname{tr}_U(BA)=\operatorname{tr}_V(AB)$: w bazach oba ślady są sumą tych samych iloczynów wpisów macierzy, także gdy wymiary się różnią. Stąd
\[
 \begin{aligned}
 T_\alpha(\mu)&=\operatorname{tr}_{L(\lambda)_{\mu+\alpha}}
                       (X_\alpha Y_\alpha)\\
 &=T_\alpha(\mu+\alpha)
       +(\mu+\alpha,\alpha)m_{\mu+\alpha}.
 \end{aligned}
\]
Iterujemy tę równość. Dla dostatecznie dużego $N$ przestrzeń wagi $\mu+N\alpha$ jest zerowa, ponieważ moduł ma skończony zbiór wag i $\alpha\ne0$ w charakterystyce zero. Otrzymujemy $T_\alpha(\mu)=\sum_{j\ge1}(\mu+j\alpha,\alpha)m_{\mu+j\alpha}$. Biorąc ślad wzoru na Casimir z poprzedniego lematu na przestrzeni wagi $\mu$, dostajemy
$c_\lambda m_\mu=(\mu,\mu+2\rho)m_\mu+2\sum_{\alpha\in\Phi^+}T_\alpha(\mu)$.
Podstawienie daje tezę. Wszystkie ślady i wymiary zerowych przestrzeni interpretujemy jako zero, więc dowód obejmuje także wagi spoza nośnika.
\end{proof}

\section{Operator różniczkowy na pierścieniu grupowym}

Pracujemy w $\mathbb Q[P]$, czyli w pierścieniu grupowym z wymiernymi współczynnikami. Dla $v\in P\otimes_{\mathbb Z}\mathbb Q$ definiujemy liniowo $\partial_v e^\mu=(v,\mu)e^\mu$ i $\mathcal L e^\mu=\|\mu\|^2 e^\mu$. Pierścień $\mathbb Q[P]$ jest dziedziną, ponieważ jest pierścieniem wielomianów Laurenta w zmiennych odpowiadających wagom fundamentalnym (AA-68). Możemy więc dzielić przez jego niezerowe elementy w ciele ułamków.

\begin{lemat}
\label{aa:lem:weyl-differential-product}
Przy powyższych założeniach dla $F\in\mathbb Q[P]$ zachodzi w ciele ułamków
\[
 \begin{aligned}
 D^{-1}\mathcal L(DF)
 &=\mathcal LF+2\partial_\rho F+\|\rho\|^2 F\\
 &\quad+2\sum_{\alpha\in\Phi^+}
     \frac{e^{-\alpha}}{1-e^{-\alpha}}\partial_\alpha F.
 \end{aligned}
\]
\end{lemat}
\begin{proof}
Każde $\partial_v$ jest derywacją, ponieważ $(v,\mu+\nu)=(v,\mu)+(v,\nu)$. Dla $F=\sum_\mu f_\mu e^\mu$ i $G=\sum_\nu g_\nu e^\nu$ rozwinięcie kwadratu normy daje
\[
 \mathcal L(GF)=G\mathcal LF+F\mathcal LG+
 2\sum_{\nu}g_\nu e^\nu\partial_\nu F.
\]
Dla ustalonego $F$ oznaczmy ostatnią sumę bez czynnika $2$ przez $B_F(G)$. Jest to derywacja względem $G$ o wartościach w pierścieniu, gdyż dla monomów $B_F(e^{\nu+\eta})=e^\nu B_F(e^\eta)+e^\eta B_F(e^\nu)$, a następnie stosujemy dwuliniowość. W szczególności $B_F(e^\rho)=e^\rho\partial_\rho F$ oraz $B_F(1-e^{-\alpha})=e^{-\alpha}\partial_\alpha F$. Z reguły iloczynu dla $D=e^\rho\prod_{\alpha>0}(1-e^{-\alpha})$ otrzymujemy
\[
 D^{-1}B_F(D)=\partial_\rho F+
     \sum_{\alpha\in\Phi^+}\frac{e^{-\alpha}}{1-e^{-\alpha}}\partial_\alpha F.
\]
Mianownik jest niezerowy na mocy AA-69, a każdy czynnik jest niezerowy, gdyż $\alpha\ne0$. Ponadto $D=A_\rho$, więc $\mathcal LD=\|\rho\|^2D$ dzięki niezmienniczości normy względem $W$. Wstawienie obu równości do reguły iloczynu kończy dowód.
\end{proof}

\section{Wzór charakterów}

\begin{twierdzenie}[Wzór charakterów Weyla]
\label{aa:tw:weyl-character-formula}
Przy powyższych założeniach dla każdej $\lambda\in P_+$ zachodzi w $\mathbb Z[P]$ tożsamość
\[
 A_\rho\operatorname{ch}L(\lambda)=A_{\lambda+\rho}.
\]
Równoważnie, w ciele ułamków pierścienia grupowego,
\[
 \operatorname{ch}L(\lambda)=
 \frac{\sum_{w\in W}\epsilon(w)e^{w(\lambda+\rho)}}
 {e^\rho\prod_{\alpha\in\Phi^+}(1-e^{-\alpha})}.
\]
\end{twierdzenie}
\begin{proof}
Niech $F=\operatorname{ch}L(\lambda)$. Dla każdego dodatniego korzenia $\alpha$ połóżmy $T_\alpha=\sum_{\mu\in P}T_\alpha(\mu)e^\mu$ z oznaczeniem śladów z dowodu Lematu~\ref{aa:lem:weyl-weight-trace}. Ta suma ma skończony nośnik, gdyż ślad na zerowej przestrzeni wynosi zero. Mnożąc równość $T_\alpha(\mu)-T_\alpha(\mu+\alpha)=(\mu+\alpha,\alpha)m_{\mu+\alpha}$ przez $e^\mu$ i porównując wszystkie współczynniki, dostajemy
\[
 (1-e^{-\alpha})T_\alpha=e^{-\alpha}\partial_\alpha F.
\]
Z równania śladów Casimira wynika w pierścieniu grupowym
$c_\lambda F=\mathcal LF+2\partial_\rho F+2\sum_{\alpha\in\Phi^+}T_\alpha$.
Wszystkie współczynniki są wymierne, więc równości uzyskane nad $K$ zachodzą w $\mathbb Q[P]$, ponieważ $\mathbb Q\to K$ jest iniekcją. Podstawienie do Lematu~\ref{aa:lem:weyl-differential-product} daje
\[
 \mathcal L(DF)=(c_\lambda+\|\rho\|^2)DF
                 =\|\lambda+\rho\|^2 DF.
\]
Nie użyliśmy nieskończonych sum formalnych ani podstawiania wartości do mianownika.

Charakter $F$ jest niezmienniczy względem $W$, a $D$ antyniezmienniczy, więc $DF$ jest antyniezmienniczy. Na mocy Twierdzenia~\ref{aa:tw:weyl-alternant-basis} zapisuje się jednoznacznie jako skończona suma
$DF=\sum_{\eta\in P_+}a_\eta A_{\eta+\rho}$ z całkowitymi $a_\eta$. Z trójkątności charakteru (Lemat~\ref{aa:lem:lie-character-triangular}) i postaci iloczynowej $D$ każda waga nośnika $DF$ jest $\preceq\lambda+\rho$, a współczynnik przy $e^{\lambda+\rho}$ wynosi jeden: aby osiągnąć tę wagę, obie nieujemne kombinacje korzeni prostych odejmowane od wag najwyższych muszą być zerowe. Ponieważ nośniki alternantów są rozłączne i $e^{\eta+\rho}$ ma współczynnik jeden w swoim alternancie, mamy $a_\lambda=1$, a $a_\eta\ne0$ pociąga $\eta\preceq\lambda$.

Każdy alternant jest wektorem własnym $\mathcal L$ o wartości własnej $\|\eta+\rho\|^2$, ponieważ $W$ zachowuje normę. Porównując współczynniki przy $e^{\eta+\rho}$ w równaniu dla $\mathcal L(DF)$, otrzymujemy dla $a_\eta\ne0$ równość $\|\lambda+\rho\|^2=\|\eta+\rho\|^2$. Gdyby $\eta\ne\lambda$, mielibyśmy $\lambda-\eta=\sum_i n_i\alpha_i$ z $n_i\ge0$ i co najmniej jednym $n_i>0$. W rzeczywistym modelu
\[
 \begin{aligned}
 \|\lambda+\rho\|^2-\|\eta+\rho\|^2
 &= (\lambda-\eta,\lambda+\eta+2\rho)\\
 &=\sum_i n_i(\alpha_i,\lambda+\eta+2\rho)>0.
 \end{aligned}
\]
Ostatnia nierówność wynika z $(\alpha_i,\xi)=\frac12\|\alpha_i\|^2\xi(h_i)$ i $(\lambda+\eta+2\rho)(h_i)\ge2$. Sprzeczność wyklucza wszystkie niższe alternanty. Pozostaje $DF=A_{\lambda+\rho}$, co daje obie postacie wzoru.
\end{proof}

\begin{przyklad}[Typ $A_1$]
\label{aa:prz:weyl-character-a1}
Niech $\lambda=m\omega$, $m\ge0$, i $x=e^\omega$. Wtedy $\rho=\omega$, a grupa Weyla składa się z identyczności i odbicia $\omega\mapsto-\omega$. Zatem
\[
 \operatorname{ch}L(m\omega)
 =\frac{x^{m+1}-x^{-m-1}}{x-x^{-1}}
 =\sum_{j=0}^m x^{m-2j}.
\]
Ostatnia równość wynika z pomnożenia sumy przez $x-x^{-1}$: wszystkie składniki pośrednie znoszą się parami. Zgadza się to z modelem z AA-68, również dla $m=0$.
\end{przyklad}

\begin{przyklad}[Typ $A_2$, reprezentacja dołączona]
\label{aa:prz:weyl-character-a2}
Dla $\mathfrak{sl}_3(K)$ waga najwyższa reprezentacji dołączonej wynosi $\alpha_1+\alpha_2=\rho$. Moduł jest nieprzywiedlny, gdyż $\mathfrak{sl}_3(K)$ jest prostą algebrą Liego. Z AA-69
$A_{2\rho}$ otrzymujemy z $A_\rho$ przez podstawienie $e^\mu\mapsto e^{2\mu}$. Dlatego
\[
 \operatorname{ch}L(\rho)=
 \frac{e^{2\rho}\prod_{\alpha\in\Phi^+}(1-e^{-2\alpha})}
      {e^\rho\prod_{\alpha\in\Phi^+}(1-e^{-\alpha})}
 =e^\rho\prod_{\alpha\in\Phi^+}(1+e^{-\alpha}).
\]
Dla $\Phi^+=\{\alpha_1,\alpha_2,\alpha_1+\alpha_2\}$ rozwinięcie ośmiu składników daje sześć korzeni z krotnością jeden i wagę zero z krotnością dwa (dwa podzbiory mają sumę $\rho$: jednoelementowy z $\alpha_1+\alpha_2$ i dwuelementowy z $\alpha_1,\alpha_2$). Zgadza się to z rozkładem korzeniowym $\mathfrak{sl}_3(K)$ i wymiarem osiem.
\end{przyklad}

Wzór jest tożsamością pierścienia grupowego po pomnożeniu przez mianownik. Bezpośrednie podstawienie $e^\mu=1$ w ilorazie daje postać nieoznaczoną. Wyprowadzenie wzoru wymiarowego wymaga osobnego argumentu granicznego.
